\documentclass[10pt]{article}
\usepackage[margin=20mm]{geometry}
\usepackage{amsmath,amssymb,mathtools,fontspec,unicode-math,xurl,hyperref}
\IfFontExistsTF{DejaVu Serif}{\setmainfont{DejaVu Serif}}{\setmainfont{Latin Modern Roman}}
\setmathfont{latinmodern-math.otf}
\setlength{\parindent}{0pt}\setlength{\parskip}{6pt}\sloppy
\title{Initial scientific review of
Quadratic Forms in 8 Prime Variables}
\author{Targeted recovery 0025: M25-0398}
\date{Exact-source initial review; independent scientific acceptance pending}
\begin{document}
\maketitle
Author. Ben Green. Frozen identity. Geometric and Functional Analysis 35 (2025), 1587–1637; DOI https://doi.org/10.1007/s00039-025-00727-9. Reviewed object. The exact 55-page arXiv:2108.10401v1 PDF, dated 23 August 2021 on its first-page margin. Source-page numbers below refer to this PDF. Equivalence to the journal version is not assumed.\par
\section*{Assessment and scope}
The manuscript develops a striking approach to quadratic forms in eight prime variables. Its main new ingredient is an averaged exponential-sum estimate obtained from the Weil representation, a bounded convolution argument, and the identification of a subgroup conjugate to \(\operatorname{SL}_2(\mathbb F_p[\Delta])\). The subsequent exclusion of invariant vectors is especially direct.\par
Major revision is required for the exact version reviewed. Two Major issues are established. First, the printed local density in (1.3) counts the zero level independently of \(N\). An explicit generic form then gives a positive asserted main term although the actual counting function is identically zero. Second, the subgroup conclusion in Appendix C is false even for exact equality in its norm hypothesis. This is a core auxiliary result explicitly invoked in Section 9. Its failure does not itself refute the intended, corrected asymptotic theorem: the actual measure there has additional structure, which a replacement lemma must use.\par
Eight Minor issues have bounded repairs or concern the scope of local statements. The full source, including all appendices and references, was read; the relevant source images were actually inspected. Independent exact calculations and selected complete proof units in primary references support this report. No author code was executed. This is completion of an initial scientific review, not independent acceptance.\par
\section*{Major issues}
\subsection*{P1. The local factor in the main theorem counts the wrong level}
Location. Theorem 1.1 and (1.3), p. 2; definition of \(\beta_\infty\), p. 3; the incompatible orthogonality calculation on p. 12.\par
For \[Q(x,y)=x^Tax+x^Tby+y^Tcy,\qquad
 \Delta=4b^{-1}ab^{-T}c-I,\] the theorem assumes \(\det(ab)\ne0\) and four distinct eigenvalues of \(\Delta\) in \(\overline{\mathbb Q}\setminus\{-1,0\}\). Its local factor is printed as \[\beta_{p,n}(N)=p^{-7n}
 \sum_{\substack{x,y\in(\mathbb Z/p^n\mathbb Z)^4\\Q(x,y)\equiv0\pmod {p^n}}}
 \Lambda_{\mathbb Z/p\mathbb Z}^{\otimes4}(x)\Lambda_{\mathbb Z/p\mathbb Z}^{\otimes4}(y).
 \tag{R1}\] Thus the congruence in the displayed definition has no \(N\). The source image confirms the zero; this is not a text-extraction ambiguity.\par
Counterexample satisfying all genericity assumptions. Put \[Q_0(x,y)=\sum_{i=1}^4(x_i-y_i)\bigl(i x_i+(i+1)y_i\bigr),
 \qquad Q=2Q_0,\qquad N=1.
 \tag{R2}\] Then \[a=2\operatorname{diag}(1,2,3,4),\quad b=2I,\quad
 c=-2\operatorname{diag}(2,3,4,5),\quad
 \Delta=\operatorname{diag}(-9,-25,-49,-81).\] Here \(\det a=384\), \(\det b=16\); all four eigenvalues are distinct and avoid \(0,-1\). Every integer value of \(Q\) is even. Consequently the left side of Theorem 1.1 is zero for every \(X\), including its von Mangoldt weights and prime powers.\par
It remains necessary to prove that the printed main term is positive; a finite residue check alone would not establish this. At the all-ones vector, \(Q_0=0\) and \[\nabla Q_0=(3,5,7,9,-3,-5,-7,-9).
 \tag{R3}\] For an odd prime \(p\), at least one of the binary blocks of \(Q_0\) has nonzero discriminant \((2i+1)^2\) modulo \(p\). Its Hessian is invertible. At any unit vector the gradient of that block cannot vanish. Hence every unit solution of \(Q_0=0\) modulo \(p\) is nonsingular, and each lifts to exactly \(p^7\) unit solutions at the next level. There is at least one such solution, by (R3). Multiplication by \(2\) does not change the zero congruence for odd \(p\). If \(C_p\) counts the unit zeros modulo \(p\), then \[\beta_p^{\rm printed}(1)
 =p^{-7}\left(\frac p{p-1}\right)^8 C_p>0
 \qquad(p\ \hbox{odd}).
 \tag{R4}\]\par
At \(p=2\), there is one unit vector modulo \(2\), and \(Q_0\) has an odd derivative there. The number of unit zeros of \(Q_0\) modulo \(2^k\) is \(2^{7(k-1)}\). For \(n\ge2\), the number of unit zeros of \(2Q_0\) modulo \(2^n\) is therefore \[2^8\,2^{7(n-2)}=2^{7n-6}.\] The weight in (R1) is \(2^8\), so \(\beta_{2,n}^{\rm printed}(1)=4\) for every \(n\ge2\), and \(\beta_2^{\rm printed}(1)=4\).\par
The infinite product is also positive. For \(p>7\), each binary form \(q_i(u,v)=(u-v)(iu+(i+1)v)\) is a split nondegenerate form. For \(r\ne0\), \[S_i(r):=\sum_{u,v\in\mathbb F_p^\times}e_p(rq_i(u,v))
 =p-\sum_u e_p(ri u^2)-\sum_v e_p(-r(i+1)v^2)+1.\] The two one-dimensional Gauss sums have modulus \(\sqrt p\), so \(|S_i(r)|\le p+2\sqrt p+1\). Orthogonality gives \[C_p=\frac{(p-1)^8}{p}
       +\frac1p\sum_{r\ne0}\prod_{i=1}^4S_i(r)
     =\frac{(p-1)^8}{p}+O(p^4).\] Together with (R4), this yields \(\beta_p^{\rm printed}(1)=1+O(p^{-3})\). All finitely many remaining factors are positive, so their product with the convergent tail is strictly positive.\par
Finally, the vector \(x_i=y_i=1/2\) is an interior zero of \(Q\) in \([0,1]^8\) with nonzero gradient. The implicit-function theorem on a small interior box gives a uniformly positive contribution to the level density at \(N/X^2=1/X^2\) for all sufficiently large \(X\). The possible critical point at the origin does not prevent existence of the density for this nondegenerate eight-variable quadratic form; its local contribution is integrable by quadratic scaling. Thus \(\beta_\infty\ge c_Q>0\) for large \(X\). The asserted main term is at least \(cX^6\), while the counting function is zero and the error is \(o(X^6)\) for any fixed \(A>0\). This contradicts the theorem with (R1).\par
Independent corroboration. Exact residue convolution, with no floating-point arithmetic, checked 18 prime-power cases. At \(p=2\), the printed densities for \(n=1,2,3,4,5\) are \(2,4,4,4,4\); the corrected target-one densities are all zero. These computations support, but do not replace, the all-prime argument above.\par
Impact and repair boundary. Replace the zero congruence in (1.3) by \(Q(x,y)\equiv N\pmod {p^n}\). The p. 12 calculation already uses \(Q-N\) and supports this interpretation. For (R2), the corrected \(2\)-adic factor vanishes, removing the contradiction. The severity is Major because the printed leading constant makes the main theorem false under all its hypotheses, even though the textual repair is short. This finding does not refute the theorem with the corrected local factors, and it is not transferred to the journal version.\par
\subsection*{P2. Symmetric near-extremizers need not be close to a subgroup}
Location. Lemma C.2, p. 52; the claim following (C.4) and Corollary C.3, p. 53; their use in Section 9, Step 2, p. 34.\par
Corollary C.3 assumes that \(\mu\) is a symmetric probability measure on a finite group and \[\|\mu*\mu\|_2\ge(1-\varepsilon)\|\mu\|_2.
 \tag{R5}\] It concludes that \(\mu\) is \(O(\varepsilon^c)\)-close in relative \(\ell^2\) norm to the uniform measure on a subgroup, with additional mass and support conclusions.\par
Exact counterexample. Let \(G=C_2=\{e,s\}\) with \(s^2=e\), and take \(\mu=\delta_s\). It is a symmetric probability measure and \[\mu*\mu=\delta_e,\qquad \|\mu*\mu\|_2=\|\mu\|_2=1.\] Thus (R5) holds for every positive \(\varepsilon\). The only subgroups are \(\{e\}\) and \(G\), and \[\|\delta_s-\delta_e\|_2^2=2,\qquad
 \|\delta_s-\mu_G\|_2^2=\frac12.
 \tag{R6}\] The latter distance divided by \(\|\mu_G\|_2\) equals \(1\). Neither distance can be \(O(\varepsilon^c)\) as \(\varepsilon\downarrow0\). The support conclusion also fails for \(H=G\) when \(\varepsilon\) is small. The same example embeds into \(\operatorname{Sp}_8(\mathbb F_p)\), for odd \(p\), by taking \(s=-I\). For any subgroup containing \(s\), its distance squared from \(\delta_s\) is \(1-1/|H|\ge1/2\); for a subgroup not containing \(s\), it is \(1+1/|H|\).\par
There is a second, independent normalization problem in Lemma C.2: its printed hypothesis is \[\|\mu*\mu\|_2\ge(1-\varepsilon)\|\mu\|_2^2.
 \tag{R7}\] For every symmetric probability measure, \((\mu*\mu)(e)=\|\mu\|_2^2\le\|\mu*\mu\|_2\), so (R7) imposes no near-equality restriction at all. For example, on \(C_3\), take \(\mu=(1/2,1/4,1/4)\). Exact calculation gives \[\|\mu\|_2^2=\frac38,\qquad
 \|\mu*\mu\|_2^2=\frac{43}{128},\qquad
 \|\mu-\mu_{C_3}\|_2^2=\frac1{24}.\] This also illustrates why merely adding positive definiteness cannot repair the squared norm in (R7). Changing (R7) to (R5), however, still leaves the counterexample (R6).\par
Where the proof fails. High multiplicative energy can identify a coset of a subgroup. Symmetry of a set does not force such a coset to contain the identity: \(\{s\}\) is already a symmetric coset. The transition from (C.4) to closeness of \(S\) to an unshifted subgroup omits precisely this obstruction. Theorem 3 of Fournier’s primary paper characterizes equality up to left and right translations; its proof explicitly obtains a coset. It does not supply the stronger unshifted conclusion quoted here.\par
Core impact and repair boundary. Section 9 uses Corollary C.3 to obtain (9.6)–(9.8), then identifies \(H=\Gamma_p\) and derives the bounded convolution estimate in Proposition 8.4. That estimate is one of the two central hypotheses of the amplification argument yielding Proposition 6.2. The erroneous auxiliary assertion therefore has a Major role.\par
The actual measure in Step 2 is an even symmetrized convolution power. Its convolution operator is positive semidefinite; \(\delta_s\) above does not have this property. Consequently the counterexample disproves the stated general lemma, not Proposition 8.4 for the actual measures. A repair must prove a correctly normalized lemma for this additional class, or retain the coset conclusion and rigorously eliminate the nonidentity coset using the extra structure. Positivity alone is not asserted here to have supplied that missing proof. The printed hypotheses and the existing appeal to Fournier do not complete it.\par
\section*{Minor issues with bounded repairs}
\subsection*{P3. The determinant factor in (1.4) is incorrect}
Location. Equation (1.4), p. 3. The displayed block factorization itself gives \[\det\begin{pmatrix}a&b/2\\b^T/2&c\end{pmatrix}
 =\det a\,
   \det\left(\frac14b^Ta^{-1}b\Delta\right)
 =2^{-8}(\det b)^2\det\Delta.
 \tag{R8}\] There are four factors of \(1/4\), not eight. The coefficient printed in (1.4) is \(2^{-16}\). For the generic form (R2), the actual determinant is \(893025\), while the printed expression is \(893025/256\). Exact rational arithmetic confirms both values. Replace \(2^{-16}\) by \(2^{-8}\). The nonvanishing equivalence used subsequently is unchanged, which bounds the impact.\par
\subsection*{P4. Distinct eigenvalues define an open condition}
Location. Genericity paragraph, p. 3. The condition that \(\Delta\) has distinct eigenvalues is called Zariski-closed. It is the failure of distinctness that is cut out by the discriminant. For instance, with \(a=\operatorname{diag}(1,2,3,t)\) and \(b=c=I\), the repeated-eigenvalue parameters are \(t=1,2,3\). The full genericity condition on this line excludes only \(0,1/4,1,2,3\), hence is a nonempty open subset which is not closed. Replace the relevant condition by the repeated-eigenvalue condition in the list of closed exceptional loci, or say that distinctness is open. The intended density conclusion survives.\par
\subsection*{P5. The first bound of Corollary 6.3 omits a modulus restriction}
Location. First statement of Corollary 6.3 and its box-counting proof, p. 18. The first statement asserts \(|S_{F_1,F_2}(r/q)|\ll X^8q^{-2}\) without \(q\le KX\). That restriction is printed only in the second statement. The proof of the first uses \((X/q+1)^8\ll X^8q^{-8}\), which requires it.\par
Fix any \(X\ge2\), take the generic form (R2), \(r=1\), and let both \(F_i\) be the indicator of the all-ones vector. Then \(S_{F_1,F_2}(1/q)=1\) for every \(q\), while \(X^8q^{-2}\to0\) as \(q\to\infty\). No constant depending only on the fixed form can make the assertion true. If simultaneous large \(X\) is preferred, take primes \(q>X^5\) and the same indicators. Add \(q\le KX\) to the first statement. All its uses in the stated minor-arc range satisfy this restriction.\par
\subsection*{P6. The integral dimensions and Fourier localization need correction}
Location. The definition of \(f_\varepsilon\), p. 13; (6.10)–(6.13), pp. 19–20.\par
On p. 13, \(f_\varepsilon\) uses \(\mathbb R^4\)-measure of a set of pairs \((u,v)\) and the discrete condition \(u,v\in[X]^4\). For a level-volume argument this must be \(\mathbb R^8\)-measure with \(u,v\in[0,X]^4\). Literally, the discrete set has zero Lebesgue measure, whereas \(I(0)=X^8>0\). The corrected continuous definition, Fourier inversion, and the scaling \(Q(Xx,Xy)=X^2Q(x,y)\) give the intended \(X^6\) factor.\par
On p. 19, \(w\) is typed as a function on \(\mathbb R^2\), is required to equal one on \([0,1]^{12}\), and is applied to \((u,v)\in\mathbb R^4\times\mathbb R^4\). The intended domain and cube are \(\mathbb R^8\) and \([0,1]^8\); the Fourier integral is over \(\mathbb R^8\) with dot products. Equation (6.11) also uses a normalized expectation where \(S\) is an unnormalized sum. With \(X=2\), \(F_1=F_2=1\), \(\eta=0\) and an integral phase, its two sides are \(256\) and \(1\). Replace the expectation by a sum.\par
The claimed Fourier \(L^1\) growth \(O(\max(1,|\lambda|)^4)\) can be retained, but the displayed two derivatives in each four-dimensional block do not justify it. Indeed \(|\xi|^{-2}\) is not integrable at infinity in \(\mathbb R^4\). Here is a complete replacement estimate. Choose fixed \(w\in C_c^\infty(\mathbb R^8)\) equal to one on the unit cube, set \(W_\lambda=w e(\lambda Q)\) and \(t=\max(1,|\lambda|)\). Then \[\begin{split}
 \|\widehat W_\lambda\|_1
 &\le
 \left(\int_{\mathbb R^8}(1+|\xi|^2/t^2)^{-6}\,d\xi\right)^{1/2}\\
 &\hspace{1cm}\cdot
 \left\|(1+|\xi|^2/t^2)^3\widehat W_\lambda\right\|_2
 \ll t^4.
 \end{split}
 \tag{R9}\] The first factor is \(O(t^4)\). By Plancherel, the second is bounded by normalized derivatives of orders at most six; each satisfies \(\|D^\alpha W_\lambda\|_2\ll t^{|\alpha|}\) on the fixed compact support. Thus it is \(O(1)\). This repairs (6.12) in the correct dimension and preserves the ensuing \(M^5\) estimate. These are bounded local repairs, not a counterexample to Proposition 3.4 after repair.\par
\subsection*{P7. The finite Fourier character and a normalization factor are missing}
Location. Proposition 7.2, p. 23; displayed endpoints of the proof of Proposition 7.4, p. 25; (B.3), p. 50.\par
The Fourier kernel is printed as \(e(x^Ty)\) instead of \(e_q(x^Ty)\). Under the manuscript’s convention \(e(t)=e^{2\pi i t}\), integer residue representatives make the printed kernel identically one. For \(q=3\) and \(f(y)=e_3(y_1)\), this operator sends a nonzero vector of normalized norm one to zero. It cannot be unitary. With the correct kernel, \[3^2\mathbb E_{y\in\mathbb F_3^4} e_3(y_1+x^Ty)
 =9\,1_{\{x=(-1,0,0,0)\}},\] which has normalized norm one. Exact character-exponent counts verify this without numerical cancellation.\par
The four internal operations on p. 25 use the correct \(e_q\), and Appendix B’s actual CRT calculation on p. 51 also produces \(e_q\). These give an unambiguous local correction. Restore \(e_q\) in the defective displays and restore \(q^2\) before the final double average defining \(T_{f_1,f_2}(1)\). For \(f_1=f_2=1_{\{0\}}\), the true \(T\) is \(q^{-6}\) and the printed double average is \(q^{-8}\). The conclusion of Proposition 7.4 then follows with the intended normalized Fourier transform. Neuhauser’s Theorems 4.1 and 4.3 supply the character kernel and an ordinary representation up to unit phases. No failure of the corrected Weil representation is claimed.\par
\subsection*{P8. The stabilization proof does not account for resets of \(\alpha\)}
Location. Proof of Proposition 9.3, pp. 33–34. When \(\beta\) jumps at \(j=T'_i\), the inequalities of Lemma 9.2 allow \(\alpha(T'_i+1)\) to increase. The next descending run must start there, not at \(\alpha(T'_i)\). The printed recurrence \(T_{i+1}=T_i+\alpha_*(T_i)\) does not bound these resets.\par
An exact sequence check isolates the failed deduction. Let \(\alpha_*(1)=1\) and \(\alpha_*(j)=10\) for \(j\ge2\). Set \(\alpha(1)=1\), \(\beta(1)=0\). For \(2\le j\le361\), put \[\alpha(j)=10-((j-2)\bmod10),\qquad
 \beta(j)=\min\left(36,1+\left\lfloor\frac{j-2}{10}\right\rfloor\right),\] and thereafter put \(\alpha(j)=1\), \(\beta(j)=36\). These satisfy positivity, \(\alpha(j)\le\alpha_*(j)\), \(0\le\beta(j)\le36\), monotonicity of \(\beta\), and monotonicity of \(\alpha\) whenever \(\beta\) is constant. The printed recurrence gives \(T_{37}=352\), but the first simultaneous adjacent equality occurs at \(j=361\). All inequalities used in this abstract counting step were checked exactly. This is not asserted to be a sequence arising from the manuscript’s actual measures.\par
A safe repair is \[S_1=\alpha_*(1),\qquad
 S_{i+1}=S_i+1+\alpha_*(S_i+1).\] There are at most 36 jumps. After any jump at \(j\le S_i\), the next run starts at \(j+1\) with height at most \(\alpha_*(S_i+1)\); positivity bounds its length by that height. After the last jump, the same bound forces an adjacent equality. This gives a finite \(S_{37}\) depending only on \(\alpha_*\); in the example it is \(397\). Increasing the large-prime threshold accordingly preserves the qualitative \(t=O(1)\) conclusion. The earlier descending-run display also needs its missing endpoint \(+1\).\par
\subsection*{P9. The coset-intersection union bound loses the size of the ambient set}
Location. Section 9, Step 3, p. 35. The displayed chain uses \[|R'|\ge 1-m|\mathbb F_p^\times\setminus R|>\frac{p-1}{2}.\] The intermediate quantity must be \((p-1)-m|\mathbb F_p^\times\setminus R|\). For \(p=101\), \(R=\mathbb F_p^\times\) and \(m=1\), the printed second inequality becomes \(1>50\). These choices are compatible with the set argument (take the original coset to be the whole group). The ordinary union bound gives \(100\), as required. With (9.10) and \(m=O(\log p)\), the corrected bound exceeds \((p-1)/2\) for large \(p\), so the generation argument survives.\par
\subsection*{P10. All five splitting types need not occur for a fixed generic form}
Location. Paragraph after the ring isomorphism on p. 43. The five possible degree partitions are correctly listed, but the suggestion that all occur as \(p\) varies needs a further arithmetic hypothesis. For (R2), the characteristic polynomial is \[(T+9)(T+25)(T+49)(T+81).\] For every \(p>7\), its four roots remain distinct and lie in \(\mathbb F_p\), so \(\mathbb F_p[\Delta]\cong\mathbb F_p^4\) throughout. None of the other four types occurs. This is a form satisfying the manuscript’s definition of genericity. Delete the occurrence claim or qualify it by a suitable Galois-group condition, such as a full symmetric Galois group for the quartic. The proof handles every splitting type separately and does not require all of them to occur.\par
\section*{Proof audit, dependencies, and boundaries}
Definitions and quantifiers. The theorem’s genericity is over the algebraic closure, not over \(\mathbb Q\) or \(\mathbb F_p\) itself. The bars are visible in the inspected source. Thus the diagonalization on p. 48 is legitimate over \(\overline{\mathbb F}_p\); no erroneous assumption that every good prime splits has been inferred. The form is fixed throughout the asymptotic notation, and the large-prime threshold may depend on it. The counterexample to P1 is one fixed form with one fixed target, so it respects these dependencies.\par
Circle method and major arcs. The arcs in Section 3 cover the circle but need not form a disjoint partition. The argument bounds the complement of the major arcs by the two minor-arc sets, so overlap is not a defect. The local phase must retain \(N\) as in the orthogonality formula. Zhao’s Section 3 was read as a complete relevant proof unit: its major-arc lemma assumes rank at least five, so an eight-variable nondegenerate form is within the rank requirement. The estimates use character orthogonality, summation by parts and Siegel–Walfisz, and permit the enlarged logarithmic ranges after adjusting constants. The half-integral block notation for an integer-valued quadratic form changes notation, not the residue-class summation argument. These observations do not repair P1 automatically; the statement must be corrected.\par
Minor arcs and amplification. The first minor-arc argument absorbs the pure quadratic terms into bounded weights and applies Cauchy–Schwarz to the invertible mixed block. In Section 6, the squarefree large-prime part of \(q\) receives the averaged gain and the complementary part is controlled through its large square divisor. The successive weakening of exponents is harmless after decreasing the positive saving exponent if necessary. The Fourier-localization repair (R9) preserves the claimed logarithmic budget. Schur orthogonality and repeated Cauchy–Schwarz in Proposition 8.2 were checked, including the normalized Hilbert-space norm versus the counting-measure norm for probability measures. A phase can be chosen as the unit phase of the matrix coefficient, with arbitrary phase at zero, avoiding any convention about arg.\par
Algebraic counting. Lemma 9.2 needs a brief clarification in addition to P8. The top-dimensional Lang–Weil coefficient in Proposition 9.1 may be zero for a general variety defined over \(\mathbb F_p\). For example, \(x^2-dy^2=0\) with a nonsquare \(d\) has geometric dimension one and only one rational point. Accordingly, positivity of the particular leading coefficient asserted for the collision varieties should be justified, or the proof should use an effective rational-point dimension obtained by decomposing Frobenius-stable strata. No collision variety satisfying this paper’s special construction has been exhibited as a counterexample here. This concern is recorded as a clarification, not a confirmed false application or an additional Major issue.\par
After a suitable replacement for P2, the remaining subgroup-identification steps have a coherent structure: intersecting dilated cosets, promoting a dense parameter set to all dilates, identifying the generated subgroup, improving \(\ell^2\) control to \(\ell^\infty\) by convolution, and taking Chinese remainders. The CRT generation argument permits words to be padded by identity generators; this deals with differing word lengths across primes.\par
Identification of \(\Gamma_p\) and invariant vectors. The DUL factorization, \(BC=\Delta\), and the conjugation by \(\tau=\operatorname{diag}(I,B)\) were checked. The scalar matrix in Lemma 10.5 simplifies to \[M_\theta(r,s)=
 \begin{pmatrix}
 1+\theta-s\theta/r&1/s-1/r\\
 (s-r)\theta(1+\theta)&1+\theta-r\theta/s
 \end{pmatrix}.\] Direct multiplication verifies the lower-left conjugation polynomial used to exclude triangular and subfield cases; exact rational examples independently corroborate it. The proof establishes the sufficiently-large-characteristic version needed by Proposition 10.1. Its stronger isolated \(p\ge5\) phrasing deserves an explicit small-characteristic justification if retained; no false small-prime case is claimed here.\par
Goursat’s lemma reduces possible couplings to isomorphisms, and the coefficient comparisons rule out identical and signed factor identifications unless the distinct irreducible factors coincide. Ribet’s lemma then uses perfection. The minimal-degree assertion is needed for \(\operatorname{SL}_2\), although Appendix A states the displayed bound for \(\mathrm{PSL}_2\). It can be supplied directly: restricting a nontrivial representation to the upper unipotent subgroup gives a nontrivial character, whose orbit under the diagonal torus has at least \((|k|-1)/2\) distinct characters. If all such characters were trivial, the normal closure of the unipotent subgroup would act trivially. This provides the needed bound without treating an \(\operatorname{SL}_2\) representation as automatically factoring through its quotient.\par
For Proposition 11.2, \(C\) is invertible and \(\Delta=BC\) is self-adjoint for \(x^TCy\). Distinct eigenvalues give an orthogonal eigenbasis over \(\overline{\mathbb F}_p\) with nonzero self-pairings. The Vandermonde system forces the common zero locus of \(x^TC\Delta^j x\), \(0\le j\le3\), to be \(\{0\}\). The analogous calculation applies to the Fourier support. A function and its finite Fourier transform cannot both be supported only at zero unless the function is zero. This checks the claimed exclusion of invariant vectors once the Fourier formulas are normalized correctly.\par
Source unit | Review coverage and result\par
pp. 1–5 | Identity, hypotheses, local factors, genericity and determinant; P1, P3, P4.\par
Sections 2–5 | Proof architecture, arc covering, major-arc dependency, first minor arcs; P1 and continuous-volume part of P6.\par
Section 6 | Unitarity, sharp pointwise scale, Fourier transfer, CRT and bad moduli; P5, P6.\par
Sections 7–8 | Weil realization, phases, norm conventions and amplification; P7.\par
Section 9 | Lang–Weil reduction, stabilization, cosets, convolution improvement and CRT; P2, P8, P9 and one unconfirmed clarification.\par
Sections 10–11 | Subgroup classification interface, polynomial comparisons, Goursat/Ribet, representation degree and invariant-vector argument; P10 and bounded clarifications.\par
Appendices A–C | Finite-group facts, twisted CRT construction, Young inequality and its near-equality application; P2 and P7.\par
\section*{Evidence, source discipline, and requested revision}
The retained source has SHA256\par
`a56ccdc25a94867ff7edf83ac5c038954aa`
`aaa6d57c896c097cf7b41c3d344d3`.\par
The complete hash is recorded without a line break in the accompanying manifests. The review read all 2,648 extracted source lines, corresponding to 55 pages. Thirty-five main-source page images were actually viewed, including an additional p. 43 image; twelve selected primary-reference images were also viewed. Rendering alone was not treated as viewing.\par
Independent calculations reside in `independent\_checks.py` and `independent\_calculations.json` in this paper’s verification directory. They use exact integers, rational arithmetic, residue-count convolution and character-exponent counts. Their execution receipt records the actual command and exit code. They do not execute supplied author software and are not a replacement for the proofs in this report.\par
The primary-reference check was deliberately specific. Zhao’s published Section 3, including the complete relevant major-arc proof, was read through the publisher’s full text. Neuhauser’s generator and representation construction, particularly Theorems 4.1–4.3, was checked in retained PDF images because its text encoding was unusable. Fournier’s Theorem 3 and its complete proof were read and the decisive coset passage viewed. No claim of complete re-review of these reference papers is made.\par
A bounded search of the arXiv, publisher and institutional records did not identify an exact-version code or supplementary package needed to check this theoretical proof. Absence from that search is not a claim of nonexistence and is not counted as a mathematical defect. The 2025 bibliographic identity is preserved; no inference of equivalence between the 2021 preprint and the journal text underlies the findings.\par
The revision should correct the target-level densities, replace the false general near-extremizer statement with a proved version adequate for Section 9, and implement the bounded repairs above. The principal mathematical limitation is explicit: the printed main formula is disproved, and the general auxiliary subgroup lemma is disproved; the intended asymptotic theorem with correct densities and an appropriate strengthened auxiliary argument is not disproved by this review. Independent scientific acceptance remains pending.\par
9 B. Green, Quadratic Forms in 8 Prime Variables, arXiv:2108.10401v1, 23 August 2021, exact retained source. https://arxiv.org/abs/2108.10401v1. L. Zhao, The quadratic form in nine prime variables, Nagoya Math. J. 223 (2016), 21–65, Section 3. https://doi.org/10.1017/nmj.2016.23. M. Neuhauser, An explicit construction of the metaplectic representation over a finite field, J. Lie Theory 12 (2002), 15–30, Theorems 4.1–4.3. https://www.heldermann-verlag.de/jlt/jlt12/NEUHAUSLAT2E.PDF. J. J. F. Fournier, Sharpness in Young’s inequality for convolution, Pacific J. Math. 72 (1977), 383–397, Theorem 3. https://msp.org/pjm/1977/72-2/pjm-v72-n2-p09-s.pdf.\par
\end{document}
